\chapter{Морфизм высоты покрытия в физическое время}
\label{ch:v10-nonequilibrium-bath-cover-height-physical-time}

Высота $h(n)=n$ задаёт порядок рождения, но для физического времени нужен
аффинный морфизм
\begin{equation}
 t_n=t_0+n\tau_{\rm birth}.
\end{equation}
На семи вершинах он представлен матрицей с двумя столбцами
$({\bf1},h)$. Действие кобграницы удаляет первый столбец и сохраняет второй:
\begin{equation}
 B_{\rm cov}^{\top}({\bf1},h)=({\bf0},{\bf1}).
\end{equation}
Поэтому начало $t_0$ является аффинной координатной свободой, тогда как
$\tau_{\rm birth}$ определяет все физические интервалы между рождениями.

Предыдущие родительские связи допускают условное отождествление
\begin{equation}
 \tau_{\rm birth}=\frac{\hbar}{E_C}
 =\frac{\ell_{\rm edge}}{c}=\Delta t.
\end{equation}
Три соответствующих инварианта
\begin{equation}
 \frac{E_C\tau_{\rm birth}}{\hbar}=1,qquad
 \frac{c\tau_{\rm birth}}{\ell_{\rm edge}}=1,qquad
 \frac{E_C\ell_{\rm edge}}{\hbar c}=1
\end{equation}
совместимы с положительным родителем, чей гессиан в пространстве инвариантов
равен $I_3$. Но третье условие является отношением первых двух и не добавляет
нового размерного ранга.

Введём логарифмические переменные
\begin{equation}
 (\tau_{\rm birth},E_C,\ell_{\rm edge},c).
\end{equation}
Размерная карта имеет ранг/ядро $2/2$. После фиксации физической скорости
$c$ остаётся ранг/ядро $3/1$ и масштабная мода
\begin{equation}
 (\tau_{\rm birth},E_C,\ell_{\rm edge},c)
 \longmapsto (s\tau_{\rm birth},E_C/s,s\ell_{\rm edge},c).
\end{equation}
Она одновременно растягивает такт и ребро, сохраняя все безразмерные фазы.
Только независимый якорь длины или энергии закрывает карту до ранга $4$.

\begin{theorem}[Условное аффинное чтение времени]
Высота универсального покрытия допускает точный типизированный морфизм
$t_n=t_0+n\tau_{\rm birth}$. Все интервалы независимы от выбора $t_0$, а
часовая и причинно-метрическая формулы дают одно и то же относительное
значение тика.
\end{theorem}

\begin{nogocriterion}[Согласованный такт не является абсолютной секундой]
Аффинная свобода $t_0$ является безвредным выбором начала отсчёта. Однако
мультипликативная свобода $\tau_{\rm birth}\mapsto s\tau_{\rm birth}$
является физическим масштабным дефицитом. Три взаимно согласованные формулы
не устраняют её, поскольку одна из них алгебраически зависима.
\end{nogocriterion}

\begin{protocol}\begin{enumerate}
\item условная архитектура морфизма времени: \textbf{$10/10$};
\item совместимость трёх относительных формул: \textbf{$3/3$};
\item выделение калибровочной свободы начала времени: \textbf{$1/1$};
\item ранг/ядро после фиксации $c$: \textbf{$3/1$};
\item физическое происхождение самого временного чтения: \textbf{$0/1$};
\item происхождение абсолютной длительности рождения: \textbf{$0/1$};
\item следующий гейт: \textbf{аудит кандидатов абсолютного масштаба такта}.
\end{enumerate}\end{protocol}

Машинный сертификат сохранён в
\path{s2t_v10_cell_birth_four_volume_nonequilibrium_bath_universal_cover_birth_height_physical_time_morphism_origin_gate_results.json}.